% !TeX root = ../../Book2.tex
% 纲要标题：## 第8章　高阶同伦、纤维化与长正合列
% 纲要位置：几何学计划纲要.md，第 1300 行。
% 正文状态：初稿；尚未正式审核。

% 纲要内容开始
% 
% > 本章目标与证明范围：以提升问题连接高阶同伦和纤维化；同时处理相对群、基本群作用及低阶集合项。覆盖只是离散纤维的特例，不将纤维丛和一般纤维化等同。
% 
% 纲要内容结束

\chapter{高阶同伦、纤维化与长正合列}
\label{g2:ch:08}
\BookChapterNumber{g2:ch:08}

\begin{chapterabstract}
	球面映射能否收缩、底空间中的变形能否提升，是本章的两个基本问题。高阶同伦群以立方体拼接组织球面映射，相对同伦群则允许一个边界面在指定子空间中移动。由此构造空间对的长正合列，并说明基点变化、基本群作用和低阶尖点集的含义。Serre 与 Hurewicz 纤维化以同伦提升性质刻画映射；路径空间、拉回及纤维化替换提供具体模型。纤维化的提升把底空间同伦群识别为总空间的相对同伦群，进而给出自然长正合列。复与四元数 Hopf 丛、经典群的球面丛将这些构造用于低阶计算，并给出同调上为零而在同伦上具有无穷阶的球面映射。
\end{chapterabstract}

\begin{learninggoals}
	\begin{itemize}
		\item 在球面、立方体与空间对模型之间转换，构造拼接运算并解释高阶交换性。
		\item 用道路追踪基点变化，辨认相对同伦序列中的群、尖点集及其作用。
		\item 写出两种同伦提升性质，构造路径纤维化、拉回与映射的纤维化替换。
		\item 由提升的顶面构造连接映射，证明长正合列的正合性、自然性及比较结论。
		\item 建立两种 Hopf 丛与经典群的局部平凡化，计算低阶同伦群，并检测非平凡球面映射。
	\end{itemize}
\end{learninggoals}

在单位三维球面 \(S^3\subset\mathbb C^2\) 上，一个点 \((z_0,z_1)\) 张成一条复直线。忘掉共同相位，就把这个点送到复射影直线 \(\mathbb{CP}^1\cong S^2\)。每条纤维是一个圆周；在小范围内可以连续选取相位，沿整个球面作这种选择却会遇到障碍。这个投影能否连续变成常值映射？同调群暂时不能回答：\(S^3\) 的三次同调为 \(\mathbb Z\)，而 \(S^2\) 的三次同调为零，所以投影在这一次数必定诱导零映射。

\ProseParagraphBreak

要直接研究投影本身的变形，我们把它看成一个三维球面在目标空间中的同伦类。一般地，\(\pi_n(X)\) 保存 \(S^n\to X\) 的带基点同伦类。它能提出同调群未能直接解决的填充问题，但运算与计算也需要新的工具。对上面的投影，如果将底空间中的球面从一面开始向上提升，最后留下的边界会落在圆周纤维中。边界的绕行恰好把二维球面映射与一维环路联系起来，成为长正合列中的连接映射。

\ProseParagraphBreak

这种提升方法依靠对参数和边界的控制。第一节先在立方体的各个面上规定哪些值固定、哪些值可以移动，从而建立相对同伦序列；第二节证明纤维丛和路径空间具有所需的提升能力；第三节把两项构造接起来，逐项解释正合性及低阶分支作用。第四节给出可直接计算的局部平凡化，并用长正合列证明两种 Hopf 映射都有无穷阶。

\ProseParagraphBreak

本章使用第1章的道路与同伦、基本群、覆盖提升及相对 CW 构造。用胞腔逼近证明低维球面到高维球面的映射可缩时，会恢复固定基点这一条件；检测球面恒等映射的整数倍时，会用第3章的映射度及其同伦不变性。矩阵群的例子还用有限维线性代数与正交化，四元数的乘法约定在使用前给出。进一步阅读可参看\cite[第4.1节；第4.2节，Fiber Bundles]{Hatcher2002}。

% !TeX root = ../../Book2.tex
% 纲要标题：### 8.1 高阶与相对同伦群
% 纲要位置：几何学计划纲要.md，第 1254 行。
% 正文状态：未撰写；本文件目前是框架。

% 纲要内容开始
% 
% * 球面、立方体和空间对模型
% * 拼接运算及高阶交换性
% * 基点变化与基本群作用
% * 相对同伦序列的低阶解释
% 
% 纲要内容结束

\section{高阶与相对同伦群}
\label{g2:sec:0801}

待撰写

\subsection{球面、立方体和空间对模型}
\label{g2:subsec:0801:01}
待撰写

\subsection{拼接运算及高阶交换性}
\label{g2:subsec:0801:02}
待撰写

\subsection{基点变化与基本群作用}
\label{g2:subsec:0801:03}
待撰写

\subsection{相对同伦序列的低阶解释}
\label{g2:subsec:0801:04}
待撰写
% 纲要条目排为编号小节标题，并进入目录与书签；说明段仍保留于纲要注释中。

% !TeX root = ../../Book2.tex
% 纲要标题：### 8.2 纤维化和替换构造
% 纲要位置：几何学计划纲要.md，第 1311 行。
% 正文状态：初稿；尚未正式审核。

% 纲要内容开始
% 
% * Serre 与 Hurewicz 提升性质
% * 路径空间和同伦纤维
% * 纤维丛的同伦提升
% * 拉回和映射的纤维化替换
% 
% 纲要内容结束

\section{纤维化和替换构造}
\label{g2:sec:0802}

设一个点在底空间 \(B\) 中移动，而我们已经在它的起点上方选定了 \(E\) 中的一点。覆盖的道路提升给出一种跟随它移动的方法。如果同时移动的是一族点，就需要提升整个同伦。这个要求能够把底空间中的零伦搬到总空间，从而把一个同伦类是否消失的问题转成纤维中的边界问题。纤维化正是用这种提升能力描述映射。

\subsection{Serre 与 Hurewicz 提升性质}
\label{g2:subsec:0802:01}

\begin{definition}[同伦提升与纤维化]\label{g2:def:serre-fibration}
	设 \(p:E\rightarrow B\) 连续，\(Z\) 为拓扑空间。若对任意连续映射 \(a:Z\rightarrow E\) 与 \(H:Z\times I\rightarrow B\)，只要 \(H(z,0)=p(a(z))\)，就存在连续映射 \(\widetilde H:Z\times I\rightarrow E\)，满足
	\[
		p\widetilde H=H,\qquad \widetilde H(z,0)=a(z),
	\]
	则称 \(p\) 对 \(Z\) 具有\term[tongluntishengxingzhi]{同伦提升性质}[homotopy lifting property]。若它对每个闭圆盘 \(D^m\)（\(m\ge0\)）都有此性质，称 \(p\) 为\term[Serrexianweihua]{Serre 纤维化}[Serre fibration]；若它对每个拓扑空间 \(Z\) 都有此性质，称 \(p\) 为\term[Hurewiczxianweihua]{Hurewicz 纤维化}[Hurewicz fibration]。\footnote{Witold Hurewicz（1904--1956），波兰数学家，后在美国工作，系统发展了高阶同伦群，并在维数理论中作出重要贡献。}这里不另外要求 \(p\) 满射。
\end{definition}

提升数据组成下图；虚线同时满足两个三角形中的等式。
\[
	\begin{tikzcd}[column sep=large,row sep=large]
		Z\times\{0\} \arrow[r,"a"] \arrow[d,hook] & E \arrow[d,"p"] \\
		Z\times I \arrow[r,"H"'] \arrow[ur,dashed,"\widetilde H"] & B
	\end{tikzcd}
\]
取 \(Z=D^0\) 得到道路提升，取 \(Z=I\) 得到一族道路的提升。Hurewicz 纤维化当然是 Serre 纤维化。后者只要求闭圆盘上的提升，却足以处理球面映射及其同伦，因为这些问题都由有限个圆盘及其边界组成。

\begin{example}[乘积投影与覆盖]\label{g2:exm:product-cover-fibrations}
	对任意拓扑空间 \(B,F\)，投影 \(p:B\times F\rightarrow B\) 是 Hurewicz 纤维化。把初始提升写成 \(a(z)=(H(z,0),v(z))\)，则
	\[
		\widetilde H(z,t)=(H(z,t),v(z))
	\]
	就是所需提升。覆盖也是 Hurewicz 纤维化：\cref{g2:thm:covering-lifts}已经对任意参数空间证明了连续同伦的提升。覆盖提升由初值唯一决定；乘积投影的提升通常有很多，例如还可以在 \(F\) 内连续移动第二个坐标。
\end{example}

在证明长正合列时，常常还要指定同伦在参数空间边界上的提升。下面的相对形式说明，这些边界条件可以一起保持。

\begin{lemma}[相对 CW 对上的提升]\label{g2:lem:relative-homotopy-lift}
	设 \(p:E\rightarrow B\) 为 Serre 纤维化，\((K,L)\) 为相对 CW 复形。给定连续映射 \(H:K\times I\rightarrow B\)，以及定义在
	\[
		A=(K\times\{0\})\cup(L\times I)
	\]
	上的连续提升 \(a:A\rightarrow E\)，满足 \(pa=H|_A\)，则 \(a\) 可以延拓成 \(H\) 在整个 \(K\times I\) 上的提升。特别地，对一个闭立方体，若已在除某一个面之外的全部边界面上给出相容提升，它就能延拓到整个立方体。
\end{lemma}
\begin{proof}
	先处理一个新胞腔。需要提升的空间是 \(C=D^m\times I\)，已经指定提升的部分为
	\[
		A_m=(D^m\times\{0\})\cup(\partial D^m\times I).
	\]
	空间对 \((C,A_m)\) 同胚于 \((D^m\times I,D^m\times\{0\})\)。可以在圆柱边界上看见这个同胚：沿每条从下底中心经过下底边缘、侧面、上底边缘而到上底中心的折线，保持角方向，把“下底加侧面”所占的参数区间线性缩为下底所占的区间，把余下区间线性伸长。这样得到边界的自同胚，将 \(A_m\) 送到下底；再沿从圆柱中心出发的径向线段延拓，就得到整个圆柱的自同胚。\(m=0\) 时两对原本相同。把给定的底映射与已知提升都通过这个同胚搬运，Serre 提升性质便给出 \(C\) 上的提升，再搬回去。这个论证只用面上的相容性，对角点附近的映射没有额外要求。

	现在从 \(L\times I\) 和整个 \(K\times\{0\}\) 出发，按相对胞腔的维数递增构造。处理一个特征映射 \(\phi:D^m\rightarrow K\) 时，下底的提升由初值给定，侧面的提升已在前一骨架上构造；两者在交集一致。刚才的圆柱结论将它们延到 \(D^m\times I\)。同一维的不同胞腔在边界上都沿用旧提升，因而能粘合。\cref{g2:lem:weak-product}保证，只要在 \(L\times I\) 及每个特征圆盘乘 \(I\) 上连续，所得映射就在 \(K\times I\) 上连续；因此胞腔可以无限多个。

	立方体的一个面可取为上底，其余面正是下底与各侧面的并。圆盘与立方体的同胚把单胞腔结论转为最后的断言；任意指定的面都可通过交换、反转坐标移到上底。
\end{proof}

\subsection{路径空间和同伦纤维}
\label{g2:subsec:0802:02}

路径本身也可以作为空间中的点。赋予 \(C(I,B)\) 紧开拓扑，其子基由“某一紧集 \(K\subset I\) 的像落在给定开集 \(U\subset B\) 内”的路径组成。在\cref{g2:lem:quotient-interval}的证明中已经建立了这里需要的指数律：对任意拓扑空间 \(Z\)，映射 \(Z\rightarrow C(I,B)\) 连续，当且仅当其伴随映射 \(Z\times I\rightarrow B\) 连续。证明的两个关键点是 \(I\) 的局部紧 Hausdorff 性和紧集的有限覆盖性质。因此，下面的路径族可以通过含一个额外区间参数的公式直接检验连续性。

\begin{definition}[基点道路空间与回路空间]\label{g2:def:path-loop-space}
	设 \((B,b_0)\) 为带基点拓扑空间。在紧开拓扑下令
	\[
		P_{b_0}B\coloneqq\{\gamma:I\rightarrow B:\gamma(0)=b_0\},\qquad
		\Omega_{b_0}B\coloneqq\{\gamma\in P_{b_0}B:\gamma(1)=b_0\}.
	\]
	两者分别称为\term[jidiandaolukongjian]{基点道路空间}[based path space]与\term[huilukongjian]{回路空间}[loop space]。它们都以常值道路为基点。记终点求值映射为 \(\operatorname{ev}_1(\gamma)=\gamma(1)\)。\symidx{\(P_{b_0}B\)：从 \(b_0\) 出发的道路空间}\symidx{\(\Omega_{b_0}B\)：以 \(b_0\) 为基点的回路空间}
\end{definition}

\begin{proposition}[道路空间的纤维化]\label{g2:prop:path-fibration}
	对任意带基点拓扑空间 \((B,b_0)\)，终点求值映射
	\[
		\operatorname{ev}_1:P_{b_0}B\longrightarrow B
	\]
	是 Hurewicz 纤维化，其在 \(b_0\) 上的纤维为 \(\Omega_{b_0}B\)，而总空间 \(P_{b_0}B\) 可缩。
\end{proposition}
\begin{proof}
	设 \(Z\) 为任意拓扑空间，初始提升给出连续道路族 \(\alpha_z\)，底同伦为 \(H:Z\times I\rightarrow B\)，满足 \(H(z,0)=\alpha_z(1)\)。移动到时刻 \(s\) 时，把原道路与已经走过的底道路接在一起；同时让后一段所占时间从零连续增长。具体定义
	\[
		\lambda_{z,s}(t)=
		\begin{cases}
			\alpha_z((1+s)t),&0\le t\le(1+s)^{-1},\\
			H(z,(1+s)t-1),&(1+s)^{-1}\le t\le1.
		\end{cases}
	\]
	在接缝处，两式均等于 \(\alpha_z(1)=H(z,0)\)。它们分别定义在两个闭集上并且连续，所以 \((z,s,t)\mapsto\lambda_{z,s}(t)\) 连续。指数律给出连续的道路族 \(\widetilde H(z,s)=\lambda_{z,s}\)。其起点恒为 \(b_0\)，终点为 \(H(z,s)\)，且 \(s=0\) 时恰为原道路 \(\alpha_z\)，故这是所需提升。

	定义收缩
	\[
		K(\gamma,s)(t)=\gamma((1-s)t).
	\]
	指数律再次保证其连续。它从恒等映射变到常值道路，始终固定常值道路，故 \(P_{b_0}B\) 可缩。纤维的描述直接来自终点条件。
\end{proof}

这个例子把“球面映射怎样移动”转成了“回路怎样成族”：总空间中的道路可以缩回起点，底空间的同伦信息于是集中在回路空间中。下一节的长正合列将把这种观察写成 \(\pi_n(B)\cong\pi_{n-1}(\Omega_{b_0}B)\)。

\begin{definition}[同伦纤维]\label{g2:def:homotopy-fiber}
	设 \(f:X\rightarrow B\) 连续，\(b_0\in B\)。赋予下列空间以 \(X\times C(I,B)\) 的子空间拓扑：
	\[
		\operatorname{hofib}_{b_0}(f)\coloneqq
		\{(x,\gamma):\gamma(0)=f(x),\ \gamma(1)=b_0\}.
	\]
	称它为 \(f\) 在 \(b_0\) 上的\term[tonglunxianwei]{同伦纤维}[homotopy fiber]。若已选 \(x_0\in X\) 且 \(f(x_0)=b_0\)，它的基点取 \((x_0,c_{b_0})\)，其中 \(c_{b_0}\) 是常值道路。\symidx{\(\operatorname{hofib}_{b_0}(f)\)：映射在 \(b_0\) 上的同伦纤维}
\end{definition}

通常纤维要求 \(f(x)=b_0\)，同伦纤维则把从 \(f(x)\) 到 \(b_0\) 的道路也保留下来。这样，底点沿道路移动时，已经选择的路径可以随之拼接。比如 \(X\) 只有一个点且 \(f(*)=b_0\) 时，同伦纤维就是 \(\Omega_{b_0}B\)。

\subsection{纤维丛的同伦提升}
\label{g2:subsec:0802:03}

局部乘积中的提升只要保持纤维坐标不变。要把这种办法用于纤维丛，需要让不同局部乘积中的选择在交界处接合。闭圆盘的紧性使我们可以把整个提升问题分成有限多个落在单个平凡化邻域内的小问题。

回顾\cref{g2:def:topological-fiber-bundle}：连续满射 \(p:E\to B\) 局部平凡，指每个 \(b\in B\) 都有开邻域 \(U\)、一个纤维模型 \(F_U\) 及同胚
\[
	h_U:p^{-1}(U)\longrightarrow U\times F_U,
	\qquad \operatorname{pr}_U\circ h_U=p|_{p^{-1}(U)}.
\]
当这些模型都可选为同一个非空空间 \(F\) 时，称其纤维为 \(F\)。以下证明用到的是这样的局部乘积图及其连续性。

\begin{theorem}[纤维丛的提升]\label{g2:thm:bundle-serre}
	任意局部平凡纤维丛 \(p:E\rightarrow B\) 都是 Serre 纤维化，并且具有\cref{g2:lem:relative-homotopy-lift}中的相对 CW 同伦提升性质。
\end{theorem}
\begin{proof}
	闭圆盘同胚于闭立方体，所以只需处理 \(H:I^m\times I\rightarrow B\) 及其初始提升。对一组平凡化邻域 \(\{U_\alpha\}\)，逆像 \(\{H^{-1}(U_\alpha)\}\) 是紧度量空间 \(I^m\times I\) 的开覆盖。取一个 Lebesgue 数，再把空间坐标与时间坐标都作足够细的有限等距细分，使每个小棱柱 \(C\times[t_j,t_{j+1}]\) 的像包含在某个 \(U_\alpha\) 内。Lebesgue 数的存在可由有限子覆盖直接证明：若不存在，就能选取直径趋零却不落在任何一个覆盖成员内的子集，选点取收敛子列后，与极限点的一个覆盖邻域矛盾。

	按时间片顺序提升。在一片内部，先提升空间网格的顶点，再提升边、二维面，直到所有 \(m\) 维小立方体。处理一个 \(r\) 维网格胞腔 \(C\) 时，在
	\[
		A_C=(C\times\{t_j\})\cup(\partial C\times[t_j,t_{j+1}])
	\]
	上已经有提升：下底来自前一时间片，侧面来自本片较低维胞腔。取包含整个 \(H(C\times[t_j,t_{j+1}])\) 的平凡化邻域 \(U_\alpha\)。把已有提升写成
	\[
		h_\alpha\bigl(\widetilde H(z,t)\bigr)
		=\bigl(H(z,t),v(z,t)\bigr),\qquad (z,t)\in A_C.
	\]
	由\cref{g2:prop:disk-hep}，经缩放和圆盘、立方体间的同胚，有回缩 \(r_C:C\times[t_j,t_{j+1}]\rightarrow A_C\)。定义新提升为
	\[
		\widetilde H(z,t)=h_\alpha^{-1}\bigl(H(z,t),v(r_C(z,t))\bigr).
	\]
	它的投影严格等于当前 \(H(z,t)\)，并在 \(A_C\) 上等于原提升。对于零维胞腔，回缩就是退回该时间片的起点。因此整个逐维构造可以开始，也不会改变已完成的边界。有限闭集粘合保证每片连续，再逐时间片粘合，得到全部同伦的提升。由\cref{g2:lem:relative-homotopy-lift}即得相对 CW 形式。
\end{proof}

覆盖对应离散纤维的局部乘积，其提升唯一性使各局部选择自动一致。一般纤维丛需要上面保持边界的延拓步骤。纤维丛的这一提升证明也可对照 \cite[第4.2节，Fiber Bundles]{Hatcher2002}。

\begin{example}[纤维化中纤维的形状可以改变]\label{g2:exm:fibration-not-bundle}
	令
	\[
		E=\{(b,v)\in\mathbb R^2:|v|\le|b|\},
		\qquad p(b,v)=b.
	\]
	这是一个 Hurewicz 纤维化。事实上，给定任意 \(H:Z\times I\rightarrow\mathbb R\) 和初始提升 \(a(z)=(H(z,0),v_0(z))\)，令
	\[
		v(z,t)=\max\{-|H(z,t)|,\min\{v_0(z),|H(z,t)|\}\}.
	\]
	最大、最小运算连续，\((H(z,t),v(z,t))\) 始终落在 \(E\)，且因 \(|v_0(z)|\le|H(z,0)|\)，在时刻零恢复给定初值。\(p^{-1}(0)\) 是一点，\(p^{-1}(b)\) 在 \(b\ne0\) 时是非退化闭区间，所以 \(p\) 在零点附近没有局部平凡化。提升性质控制连续变形，局部平凡化还要求邻近纤维具有相同的拓扑形状。
\end{example}

\subsection{拉回和映射的纤维化替换}
\label{g2:subsec:0802:04}

先看怎样限制底空间。若 \(u:C\rightarrow B\) 连续，将 \(E\) 中位于 \(u(c)\) 上方的点与 \(c\) 配对，得到
\[
	u^*E\coloneqq\{(c,e)\in C\times E:u(c)=p(e)\},
	\qquad p_C(c,e)=c.
\]
其中取乘积空间的子空间拓扑。这个\term[lahuixianweihua]{拉回}[pullback]给出交换图
\[
	\begin{tikzcd}[column sep=large,row sep=large]
		u^*E \arrow[r,"\operatorname{pr}_E"] \arrow[d,"p_C"'] & E \arrow[d,"p"] \\
		C \arrow[r,"u"'] & B.
	\end{tikzcd}
\]
图中左上空间只由“右边投影与下边映射相等”的条件确定。

\begin{proposition}[拉回保持提升性质]\label{g2:prop:fibration-pullback}
	设 \(p:E\rightarrow B\) 与 \(u:C\rightarrow B\) 连续。若 \(p\) 为 Serre 纤维化，则 \(p_C:u^*E\rightarrow C\) 也是；若 \(p\) 为 Hurewicz 纤维化，则 \(p_C\) 也为 Hurewicz 纤维化。若 \(p\) 是以 \(F\) 为纤维的纤维丛，则 \(p_C\) 也是以 \(F\) 为纤维的纤维丛。
\end{proposition}
\begin{proof}
	设 \(H:Z\times I\rightarrow C\) 及初始提升 \(a(z)=(H(z,0),e(z))\) 已给定。对底同伦 \(uH\)，用 \(p\) 的提升性质取以 \(e(z)\) 为初值的提升 \(\widetilde H_E\)。则
	\[
		\widetilde H_C(z,t)=\bigl(H(z,t),\widetilde H_E(z,t)\bigr)
	\]
	落在 \(u^*E\)，且满足所需初值与投影等式。允许的参数空间 \(Z\) 与原提升性质相同，因此两种纤维化均得到保持。对纤维丛，在 \(u^{-1}(U)\) 上用原平凡化的第二个坐标，就得到 \(p_C^{-1}(u^{-1}(U))\cong u^{-1}(U)\times F\)。
\end{proof}

任意映射也可以加入一条路径，变成具有提升性质的映射。加入的路径可以随底同伦延长，而把它缩回初点又能恢复原来的源空间。

\begin{proposition}[映射的纤维化替换]\label{g2:prop:fibration-replacement}
	设 \(f:X\rightarrow B\) 为拓扑空间之间的连续映射。令
	\[
		E_f\coloneqq\{(x,\gamma)\in X\times C(I,B):\gamma(0)=f(x)\},
		\qquad p_f(x,\gamma)=\gamma(1),
	\]
	并取 \(j(x)=(x,c_{f(x)})\)、\(q(x,\gamma)=x\)。则 \(p_f\) 为 Hurewicz 纤维化，\(p_fj=f\)，\(qj=\operatorname{id}_X\)，且 \(jq\) 与 \(\operatorname{id}_{E_f}\) 相对 \(j(X)\) 同伦。因此 \(j:X\rightarrow E_f\) 是同伦等价。\(p_f\) 在 \(b_0\in B\) 上的实际纤维就是 \(\operatorname{hofib}_{b_0}(f)\)。\symidx{\(E_f\)：映射的道路空间替换}
\end{proposition}
\begin{proof}
	给定任意参数空间 \(Z\)、初始提升 \(a(z)=(x(z),\alpha_z)\) 及底同伦 \(H:Z\times I\rightarrow B\)，其中 \(H(z,0)=\alpha_z(1)\)。保持 \(x(z)\) 不动，并采用\cref{g2:prop:path-fibration}中的道路 \(\lambda_{z,s}\)，定义
	\[
		\widetilde H(z,s)=\bigl(x(z),\lambda_{z,s}\bigr).
	\]
	每条 \(\lambda_{z,s}\) 的起点仍是 \(\alpha_z(0)=f(x(z))\)，所以它落在 \(E_f\)；终点与初值也分别给出 \(p_f\widetilde H=H\) 和 \(\widetilde H(z,0)=a(z)\)。上述公式的连续性已由指数律证明，故 \(p_f\) 为 Hurewicz 纤维化。

	公式
	\[
		K((x,\gamma),s)=\bigl(x,t\mapsto\gamma((1-s)t)\bigr)
	\]
	给出 \(\operatorname{id}_{E_f}\) 到 \(jq\) 的同伦，并固定所有常值道路；\(qj=\operatorname{id}_X\) 与 \(p_fj=f\) 可直接求值得到。最后，\(p_f^{-1}(b_0)\) 的两个路径端点条件与同伦纤维的定义完全相同。
\end{proof}

纤维化替换把 \(f\) 分解为
\[
	X\xrightarrow[\simeq]{\ j\ }E_f\xrightarrow{\ p_f\ }B.
\]
于是可以对 \(p_f\) 使用长正合列，并用同伦等价 \(j\) 将总空间中的同伦群换回 \(X\)。这正是同伦纤维进入一般映射问题的用途。

% !TeX root = ../../Book2.tex
% 纲要标题：### 8.3 纤维化的同伦长正合列
% 纲要位置：几何学计划纲要.md，第 1318 行。
% 正文状态：初稿；尚未正式审核。

% 纲要内容开始
% 
% * 连接映射的提升构造
% * 逐项验证正合性
% * 基点、连通分支和群作用
% * 自然性与比较
% 
% 纲要内容结束

\section{纤维化的同伦长正合列}
\label{g2:sec:0803}

设 \(p:E\to B\) 是 Serre 纤维化，选定 \(e_0\in E\)，记 \(b_0=p(e_0)\)、\(F=p^{-1}(b_0)\)。底空间中的一个球面映射未必能提升为总空间中的球面映射：沿参数逐步提升以后，最后一个面可能停在 \(F\) 中的一张非平凡球面上。连接映射记录的正是这个未能收拢的面。先允许这一个面在纤维中移动，提升就总能完成；这把纤维化问题与已经建立的相对同伦联系起来。

\subsection{连接映射的提升构造}
\label{g2:subsec:0803:01}

沿用\cref{g2:def:relative-homotopy}的立方体模型，令 \(I^n=I^{n-1}\times I\)，最后一个坐标记为 \(t\)，并令
\[
	J^{n-1}=(\partial I^{n-1}\times I)\cup(I^{n-1}\times\{0\}).
\]
顶面 \(t=1\) 可以映入 \(F\)，其余面必须映到 \(e_0\)。投影到 \(B\) 后，整个边界都成为 \(b_0\)，于是有映射
\[
	P_n:\pi_n(E,F,e_0)\longrightarrow\pi_n(B,b_0),\qquad [v]\longmapsto[pv].
\]

\begin{proposition}
	\label{g2:prop:fibration-relative-comparison}
	若 \(p:E\to B\) 是 Serre 纤维化，\(e_0\in E\)、\(b_0=p(e_0)\)、\(F=p^{-1}(b_0)\)，则 \(P_n\) 在 \(n\ge2\) 时是群同构，在 \(n=1\) 时是保持基点的集合双射。
\end{proposition}

\begin{proof}
	给定 \(u:(I^n,\partial I^n)\to(B,b_0)\)，先在 \(J^{n-1}\) 上指定常值提升 \(e_0\)。\cref{g2:lem:relative-homotopy-lift}把这个提升延拓到整个立方体，得到 \(\widetilde u\)。因为 \(u\) 在顶面也等于 \(b_0\)，\(\widetilde u\) 的顶面落在 \(F\)，所以它定义相对类，且 \(P_n[\widetilde u]=[u]\)。这证明满射。

	再设 \(v_0,v_1\) 的投影由保持边界基点的同伦 \(H:I^n\times I_s\to B\) 连接。在 \(s=0,1\) 上分别指定提升 \(v_0,v_1\)，在 \(J^{n-1}\times I_s\) 上指定 \(e_0\)。它们在交叠处一致。把 \(t\) 看成提升的时间坐标，其余 \(n\) 个坐标为 \((x,s)\)，这些已知数据恰好位于底面和全部侧面；唯有 \(t=1\) 的面尚未指定。再次应用相对提升引理，得到 \(\widetilde H\)。在顶面上，\(p\widetilde H=H=b_0\)，所以顶面始终留在 \(F\)；这正是 \(v_0\) 与 \(v_1\) 的相对同伦。这个论证也适用于 \(n=1\)，直接证明了尖点集上的单射。

	当 \(n\ge2\) 时，投影保持沿第一坐标的拼接，故 \(P_n\) 是群同态。已经得到的双射因而是群同构。
\end{proof}

\begin{definition}
	\label{g2:def:fibration-connecting-map}
	在上述条件下，纤维化的\term[lian jie ying she]{连接映射}[connecting map]定义为
	\[
		\partial_n=\partial\circ P_n^{-1}:\pi_n(B,b_0)\longrightarrow\pi_{n-1}(F,e_0).
	\]
	这里右侧的 \(\partial\) 是空间对 \((E,F)\) 的顶面限制。因此，对上述提升 \(\widetilde u\)，有
	\[
		\partial_n[u]=[\widetilde u(-,1)].
	\]
	当 \(n=1\) 时，右侧指提升道路终点所在的道路连通分支。
\end{definition}

\cref{g2:prop:fibration-relative-comparison}同时解释了构造为何与选择无关：同一底空间同伦类的所有提升给出同一相对类，顶面因而给出同一纤维同伦类。\(n\ge2\) 时连接映射是群同态；\(n=1\) 时其值域为 \(\pi_0(F)\)，通常只具有带基点集合的结构。

\subsection{逐项验证正合性}
\label{g2:subsec:0803:02}

连接映射的意义可以表述为一个可检验的问题：顶面能否在纤维中缩掉？若能，就可以把提升补成一个球面。反过来，如果已经有一个总空间中的球面，它的投影便不会留下这个障碍。下面把这两个方向与另外两处核、像关系一起证明。

\begin{theorem}
	\label{g2:thm:fibration-homotopy-sequence}
	设 \(p:E\to B\) 是 Serre 纤维化，\(e_0\in E\)，\(b_0=p(e_0)\)，\(F=p^{-1}(b_0)\)，\(i:F\hookrightarrow E\) 为包含。存在自然的长正合列
	\[
		\cdots\longrightarrow\pi_n(F,e_0)\xrightarrow{i_*}\pi_n(E,e_0)
		\xrightarrow{p_*}\pi_n(B,b_0)\xrightarrow{\partial_n}\pi_{n-1}(F,e_0)
		\longrightarrow\cdots
	\]
	其低阶部分为
	\[
		\begin{aligned}
			\pi_2(B,b_0)&\xrightarrow{\partial_2}\pi_1(F,e_0)
			\xrightarrow{i_*}\pi_1(E,e_0)\xrightarrow{p_*}\pi_1(B,b_0)\\
			&\xrightarrow{\partial_1}\pi_0(F,e_0)\xrightarrow{i_*}\pi_0(E,e_0)
			\xrightarrow{p_*}\pi_0(B,b_0).
		\end{aligned}
	\]
	群项上的正合指像等于核；尖点集项上的正合指前一映射的像等于后一映射所取基点的原像。若 \(B\) 道路连通，则 \(\pi_0(F)\to\pi_0(E)\) 满射，可以在其后接一个单点集。
\end{theorem}

\begin{proof}
	\cref{g2:thm:pair-homotopy-sequence}给出 \((E,F,e_0)\) 的相对同伦序列。以\cref{g2:prop:fibration-relative-comparison}中的 \(P_n\) 替换相对项，就得到定理中的映射：\(\pi_n(E)\) 先相对化再由 \(P_n\) 投影，正是通常的 \(p_*\)。下面也直接辨认每一种正合位置，以说明各个核代表什么几何事实。

	在 \(\pi_n(E)\) 处，纤维中的映射投影为常值，所以 \(\operatorname{im}i_*\subseteq\ker p_*\)。若带基点球面映射 \(a\) 的投影有带基点零伦，则用球面及其基点组成的 CW 对，提升这个零伦并令基点始终为 \(e_0\)。初始提升为 \(a\)，最终提升整个落在 \(F\)。因此 \([a]\) 来自 \(\pi_n(F)\)。这一论证对 \(n\ge1\) 都成立。

	在 \(\pi_n(B)\) 处，\(n\ge2\) 时，先取 \(u\) 的提升 \(v\)，其顶面为 \(a\)。若 \([a]=0\in\pi_{n-1}(F)\)，取 \(a\) 在 \(F\) 中的带基点零伦 \(A(x,s)\)。在 \(v\) 顶端接上这个零伦，得到
	\[
		w(x,t)=
		\begin{cases}
			v(x,2t),&0\le t\le\tfrac12,\\
			A(x,2t-1),&\tfrac12\le t\le1.
		\end{cases}
	\]
	\(w\) 的整个边界都等于 \(e_0\)，而 \(pw\) 等于先运行 \(u\)、再停在 \(b_0\) 的重参数映射，故 \([u]=p_*[w]\)。若 \([u]\) 一开始来自 \(\pi_n(E)\)，便可选整个边界为 \(e_0\) 的提升，连接类为零。\(n=1\) 时，同一个构造把“顶面零伦”换成“在 \(F\) 内从提升终点走回 \(e_0\) 的道路”，证明 \(\partial_1^{-1}([e_0])=\operatorname{im}p_*\)。

	在 \(\pi_{n-1}(F)\) 处，\(n\ge2\) 时，连接映射给出的顶面连同提升立方体，已经提供它在 \(E\) 中的零伦。反过来，若 \(a:I^{n-1}/\partial I^{n-1}\to F\) 在 \(E\) 中有带基点零伦，反向运行这个零伦就得到 \(v:I^n\to E\)，底面和侧面为 \(e_0\)，顶面为 \(a\)。此时 \(pv\) 的整个边界为 \(b_0\)，而 \(\partial_n[pv]=[a]\)。这也包括可能非交换的 \(\pi_1(F)\)。

	在 \(\pi_0(F)\) 处，连接映射得到的点由提升道路与 \(e_0\) 相连，所以它落在 \(\pi_0(E)\) 的基点原像中。反过来，若 \(y\in F\) 在 \(E\) 中与 \(e_0\) 由道路 \(c\) 相连，则 \(pc\) 是 \(b_0\) 处的环路，且 \(\partial_1[pc]=[y]\)。在 \(\pi_0(E)\) 处，若 \(e\in E\) 的投影在 \(b_0\) 的分支中，提升从 \(p(e)\) 到 \(b_0\) 的道路，得到从 \(e\) 到 \(F\) 的道路；反方向由 \(p(F)=\{b_0\}\) 立即得到。若 \(B\) 道路连通，这就对 \(E\) 的所有分支成立。
\end{proof}

\begin{example}
	\label{g2:exm:loop-space-shift}
	令 \(B\) 是带基点空间，\(P_{b_0}B\) 为从 \(b_0\) 出发的道路空间，终点映射的纤维为回路空间 \(\Omega B\)。\cref{g2:prop:path-fibration}给出这个 Hurewicz 纤维化。\(P_{b_0}B\) 由 \(\gamma(t)\mapsto\gamma((1-s)t)\) 保持常值道路地收缩，所以长正合列给出
	\[
		\pi_n(B,b_0)\cong\pi_{n-1}(\Omega B,c_{b_0})\quad(n\ge2),
		\qquad \pi_1(B,b_0)\cong\pi_0(\Omega B).
	\]
	后一个式子的集合双射在回路拼接下也是群同构。对 \(u(x,t)\)，可以取提升 \(\widetilde u(x,t)(r)=u(x,tr)\)；其顶面正是 \(x\mapsto(t\mapsto u(x,t))\)。这说明指标降低一次来自把一个参数变成回路参数。
\end{example}

\subsection{基点、连通分支和群作用}
\label{g2:subsec:0803:03}

序列末端还能记录总空间有多少个道路连通分支。沿底空间中的环路提升，一个纤维分支可能被带到另一个分支；能够这样互相到达的分支，在总空间中恰好连在一起。这是第1章覆盖单值作用在一般纤维化中的延续。

\begin{proposition}
	\label{g2:prop:fibration-component-action}
	设 \(p:E\to B\) 是 Serre 纤维化，\(B\) 道路连通，且 \(F=p^{-1}(b_0)\ne\varnothing\)。沿环路提升的终点定义右作用
	\[
		\pi_0(F)\times\pi_1(B,b_0)\longrightarrow\pi_0(F),\qquad
		[y]\cdot[\beta]=[\widetilde\beta(1)],\quad\widetilde\beta(0)=y.
	\]
	此作用的轨道集合与 \(\pi_0(E)\) 自然对应。对 \(e_0\in F\)，连接映射为 \(\partial_1[\beta]=[e_0]\cdot[\beta]\)，而 \([e_0]\) 的稳定子是 \(p_*\pi_1(E,e_0)\)。
\end{proposition}

\begin{proof}
	固定一个起点，两条提升的终点在 \(F\) 中相连：把两条提升放在一个方形的两条相对边上，在初始边指定该起点，底空间映射取与第二参数无关的原道路；缺少终点边的三条边组成提升引理中的部分边界。补全提升后，终点边就是所需的纤维内道路。若起点在 \(F\) 中沿一条道路变化，把这条道路作为方形的初始边，同样得到终点分支不变。底空间道路的端点固定同伦也由相同的方形提升处理。因此作用良定义。

	道路乘法 \(\alpha*\beta\) 先走 \(\alpha\)，后走 \(\beta\)。依次提升两段，就有 \(([y]\cdot[\alpha])\cdot[\beta]=[y]\cdot[\alpha*\beta]\)；常值道路作用为恒等，反向道路给出逆变换，所以这是右作用。作用同轨的点通过提升道路在 \(E\) 中相连；若两点在 \(E\) 中相连，投影道路就是实现相应运输的底空间环路。再由 \(B\) 道路连通时每个总空间分支都遇到 \(F\)，得到轨道集合的断言。

	若环路是 \(E\) 中基点环路的投影，就固定 \([e_0]\)。反过来，若其提升终点与 \(e_0\) 在 \(F\) 中相连，接上这条纤维内道路就得到 \(E\) 中的环路，投影与原环路同伦。因而稳定子恰为 \(p_*\pi_1(E,e_0)\)。
\end{proof}

\begin{example}
	\label{g2:exm:covering-sequence-tail}
	对 \(p:S^1\to S^1\)、\(p(z)=z^m\)，\(m\ge2\)，取基点 \(1\)。纤维是 \(m\) 次单位根，低阶序列为
	\[
		0\longrightarrow\mathbb Z\xrightarrow{\times m}\mathbb Z
		\xrightarrow{\partial_1}\{1,\zeta,\ldots,\zeta^{m-1}\}\longrightarrow\{*\},
		\qquad\zeta=e^{2\pi\mathrm i/m},
	\]
	其中 \(\partial_1(k)=\zeta^k\)。整数绕行数通过提升终点记录其模 \(m\) 的余数；纤维分支只有一个轨道，所以总空间连通。在一般连通覆盖中，纤维仍是一个传递的 \(\pi_1(B)\)-集合，稳定子可以不是正规子群。末端的集合因而要由作用理解，不能一概当作商群。
\end{example}

高阶群还需要追踪基点。若 \(\lambda\) 是 \(E\) 中从 \(e_0\) 到 \(e_1\) 的道路，\(\beta=p\lambda\)，则以 \(\lambda\) 规定移动的基点，可以运输纤维的带基点球面：对 \(a:S^r\to p^{-1}(\beta(0))\)，提升底空间同伦 \((z,t)\mapsto\beta(t)\)，同时规定 \((*,t)\mapsto\lambda(t)\)。相对提升引理给出终端球面，并由带第二同伦参数的同一构造保证其类与选择无关。反向道路提供逆映射。若 \(\lambda\) 是环路，依次提升对应于道路拼接，便得到 \(\pi_1(E,e_0)\) 在 \(\pi_r(F,e_0)\) 上的右作用；当 \(\lambda\) 位于 \(F\) 内时，它就是\cref{g2:prop:basepoint-action}中的基点变化作用。

\ProseParagraphBreak

这也解释了为何 \(\pi_0(F)\) 的运输只需底空间环路，而带基点高阶群还要指定基点的移动。选择不同的纤维内回程，最终识别之间相差一个基本群作用。实际计算中，若纤维道路连通且单连通，这个差异消失；对一般纤维，则把移动基点的道路连同群写清楚。沿 \(\lambda\) 改变总空间基点、沿 \(p\lambda\) 改变底空间基点，再用上述运输改变纤维基点，得到的长正合列彼此对应；连接映射的相容性来自对同一立方体及其顶面同时进行提升。

\subsection{自然性与比较}
\label{g2:subsec:0803:04}

纤维化经常通过拉回、坐标变换或映射路径构造与另一个纤维化比较。要把一次计算搬到新模型上，必须确认连接映射也随这些映射一起变化。它的提升定义恰好使这种相容性成为直接的复合关系。

\begin{proposition}
	\label{g2:prop:fibration-sequence-naturality}
	设 \(p:E\to B\)、\(p':E'\to B'\) 为 Serre 纤维化，带基点映射 \(h:E\to E'\)、\(g:B\to B'\) 满足 \(p'h=gp\)。记对应基点纤维为 \(F,F'\)，\(h_F=h|_F\)。则 \(h_F,h,g\) 诱导长正合列之间的态射，特别地
	\[
		\begin{tikzcd}[column sep=large]
			\pi_n(E)\arrow[r,"p_*"]\arrow[d,"h_*"']&\pi_n(B)\arrow[r,"\partial_n"]\arrow[d,"g_*"]&\pi_{n-1}(F)\arrow[d,"(h_F)_*"]\\
			\pi_n(E')\arrow[r,"p'_*"]&\pi_n(B')\arrow[r,"\partial'_n"]&\pi_{n-1}(F')
		\end{tikzcd}
	\]
	交换。\(n=1\) 时右端按带基点集合解释；\(\pi_0(F)\) 的右作用也与 \(g_*\) 和 \((h_F)_*\) 相容。
\end{proposition}

\begin{proof}
	若 \(v\) 是 \(u\) 的提升，则 \(hv\) 是 \(gu\) 的提升，且 \((hv)(-,1)=h_F(v(-,1))\)。这给出连接映射所在方块的交换性。其他方块只涉及包含与投影的复合；道路提升的终点也由 \(h_F\) 送到相应终点，故分支作用相容。
\end{proof}

下面的比较结论在替换空间或替换纤维时尤其有用。若一个连续映射在道路分支上为双射，并在每个基点的所有正次同伦群上为同构，就称它为\term[ruo tong lun deng jia]{弱同伦等价}[weak homotopy equivalence]。

\begin{proposition}
	\label{g2:prop:fibration-comparison}
	设 \(p:E\to B\)、\(p':E'\to B'\) 是 Serre 纤维化，\(B,B'\) 道路连通，连续映射 \(h:E\to E'\)、\(g:B\to B'\) 满足 \(p'h=gp\)。选 \(b_0\in B\)，令 \(F=p^{-1}(b_0)\)、\(F'=(p')^{-1}(g(b_0))\)，并假设 \(F,F'\) 非空。如果 \(g\) 与 \(h_F=h|_F:F\to F'\) 都是弱同伦等价，则 \(h\) 也是弱同伦等价。
\end{proposition}

\begin{proof}
	由\cref{g2:prop:fibration-component-action}，总空间的分支是纤维分支的作用轨道。\(\pi_1(B)\to\pi_1(B')\) 为同构，\(\pi_0(F)\to\pi_0(F')\) 为等变双射，故 \(\pi_0(E)\to\pi_0(E')\) 为双射。每个分支都含有纤维点；基点变化同构允许在这样的点上检查其余同伦群。

	固定 \(e\in F\)。先证明 \(h_*:\pi_n(E,e)\to\pi_n(E',h(e))\) 单射，其中 \(n\ge1\)。若 \(h_*x=0\)，由底空间同伦群上的单射知 \(p_*x=0\)，于是 \(x=i_*a\)。又有 \(i'_*(h_F)_*a=0\)，正合性给出 \((h_F)_*a=\partial'_{n+1}b'\)。由底空间同伦群上的满射，取 \(b\) 使 \(g_*b=b'\)。自然性和 \((h_F)_*\) 的单射给出 \(a=\partial_{n+1}b\)，故 \(x=0\)。在 \(n=1\) 时，零元按群的单位元理解，这个论证仍逐项成立。

	再给定 \(x'\in\pi_n(E')\)。取 \(b\in\pi_n(B)\) 使 \(g_*b=p'_*x'\)。\(n\ge2\) 时，\((h_F)_*\partial_n b=0\)，由纤维群上的单射知 \(\partial_n b=0\)，所以 \(b=p_*x_0\)。\(n=1\) 时用 \(\pi_0(F)\to\pi_0(F')\) 的单射得到同样结论：\(\partial_1b\) 必为基点分支，故也有 \(b=p_*x_0\)。于是 \(x'\) 与 \(h_*x_0\) 的差属于 \(\ker p'_*\)，来自某个纤维类 \(a'\)；选 \(a\) 使 \((h_F)_*a=a'\)，便得到 \(x'\) 的原像。对于 \(n=1\)，具体取 \(a'\) 满足 \(i'_*a'=x'(h_*x_0)^{-1}\)，原像为 \(i_*a\,x_0\)。这证明所有正次同伦群上的满射。
\end{proof}

\cref{g2:prop:fibration-replacement}还把任意带基点映射 \(f:X\to Y\) 换成 \(E_f\to Y\)，其中 \(E_f\simeq X\)。把这个同伦等价诱导的同构代入长正合列，得到
\[
	\cdots\longrightarrow\pi_n(F_f)\longrightarrow\pi_n(X)
	\xrightarrow{f_*}\pi_n(Y)\longrightarrow\pi_{n-1}(F_f)\longrightarrow\cdots.
\]
这里 \(F_f=\operatorname{hofib}_{y_0}(f)\)，其元素同时保存一个 \(x\in X\) 和一条从 \(f(x)\) 到 \(y_0\) 的道路。因而这个序列测量的是映射 \(f\) 的同伦纤维；即使普通原像 \(f^{-1}(y_0)\) 很小，这些道路数据仍能记录 \(f\) 对同伦群的作用。

% !TeX root = ../../Book2.tex
% 纲要标题：### 8.4 Hopf 纤维化与经典群
% 纲要位置：几何学计划纲要.md，第 1275 行。
% 正文状态：未撰写；本文件目前是框架。

% 纲要内容开始
% 
% * 复、四元数 Hopf 映射
% * 球面丛和经典群的纤维序列
% * 低阶同伦群计算
% * 由纤维化构造非平凡球面映射
% 
% ---
% 
% 纲要内容结束

\section{Hopf 纤维化与经典群}
\label{g2:sec:0804}

待撰写

\subsection{复、四元数 Hopf 映射}
\label{g2:subsec:0804:01}
待撰写

\subsection{球面丛和经典群的纤维序列}
\label{g2:subsec:0804:02}
待撰写

\subsection{低阶同伦群计算}
\label{g2:subsec:0804:03}
待撰写

\subsection{由纤维化构造非平凡球面映射}
\label{g2:subsec:0804:04}
待撰写
% 纲要条目排为编号小节标题，并进入目录与书签；说明段仍保留于纲要注释中。


\begin{chaptersummary}
	高阶同伦群把球面映射及其填充问题组织成群；两个独立的拼接方向产生交换性，相对模型则把指定顶面留给子空间。基点运输沿道路给出同构，其环路版本形成基本群作用。纤维化能够保持已有边界数据地提升同伦，因此 \(\pi_n(E,F)\) 与 \(\pi_n(B)\) 对应，顶面限制给出连接映射。长正合列的每一处核、像等式分别表达提升、填充或纤维内收缩；低阶末端通过分支与作用轨道描述总空间的连通性。路径替换使任意映射都能进入这一计算框架。Hopf 丛给出 \(\pi_2(S^2)\cong\mathbb Z\) 及球面之间的高阶比较，经典群的球面丛给出一、二次同伦群的逐阶计算；映射度对同伦加法的相容性进一步保证复与四元数 Hopf 映射具有无穷阶。
\end{chaptersummary}

\BookExercises

\begin{exercise}\label{g2:ex:homotopy-product-retract}
	设 \((X,x_0)\)、\((Y,y_0)\) 为带基点空间，\(n\ge1\)。证明两个投影诱导 \(\pi_n(X\times Y,(x_0,y_0))\cong\pi_n(X,x_0)\times\pi_n(Y,y_0)\)。由此说明 \(S^m\) 的恒等映射作为 \(S^m\times S^1\) 中的球面仍非零伦，其中 \(m\ge1\)。
\end{exercise}
\begin{solution}
	映射 \(I^n\to X\times Y\) 与一对映射 \((f_X,f_Y)\) 一一对应，固定边界条件和定边界同伦也逐分量对应。因此投影在同伦类上给出双射，逆映射为 \(([f_X],[f_Y])\mapsto[(f_X,f_Y)]\)。拼接的两段公式逐分量成立，所以这是群同构。

	令 \(i:S^m\to S^m\times S^1\)、\(i(z)=(z,1)\)，第一投影满足 \(p_1i=\operatorname{id}_{S^m}\)。若 \(i\) 零伦，复合 \(p_1\) 便使球面恒等映射零伦。但恒等映射在 \(H_m(S^m;\mathbb Z)\cong\mathbb Z\) 上诱导恒等，常值映射诱导零映射，与同调的同伦不变性矛盾。因此 \([i]\ne0\)，也相应证明恒等球面代表非零类。
\end{solution}

\begin{exercise}\label{g2:ex:relative-disk-circle}
	以 \(1\in S^1\subset D^2\) 为基点，计算 \(\pi_2(D^2,S^1,1)\) 和 \(\pi_1(D^2,S^1,1)\)。说明边界绕数为 \(k\) 的相对圆盘类如何构造，并判断两个这样的类何时相同。
\end{exercise}
\begin{solution}
	凸圆盘以 \((1-s)z+s\) 收缩到 \(1\)，且始终固定 \(1\)，故它的正次同伦群全为零。空间对序列给出
	\[
		0\longrightarrow\pi_2(D^2,S^1,1)\xrightarrow{\partial}\pi_1(S^1,1)\longrightarrow0.
	\]
	因此 \(\partial\) 是群同构，绕数进一步把它识别为 \(\mathbb Z\)。对于 \(g_k(u)=e^{2\pi\mathrm{i}ku}\)，取
	\[
		f_k(u,t)=(1-t)+t\,g_k(u),\qquad (u,t)\in I^2.
	\]
	凸性保证它进入 \(D^2\)；底边和两条侧边都为 \(1\)，顶边为 \(g_k\)，故 \(\partial[f_k]=k\)。边界同构说明每个相对类等于唯一的 \([f_k]\)，两类相同当且仅当绕数相同。

	低阶序列中 \(\pi_1(D^2)=0\)，而 \(S^1\) 道路连通。每条相对道路的终点都能在 \(S^1\) 内回到 \(1\)，从而相对同伦到圆盘中的环路；后者可收缩。因此 \(\pi_1(D^2,S^1,1)\) 只有指定元素。
\end{solution}

\begin{exercise}\label{g2:ex:relative-circle-two-points}
	令 \(X=S^1\)、\(A=\{1,-1\}\)，基点为 \(1\)。用指数覆盖计算尖点集 \(\pi_1(X,A,1)\)，并写出 \(\pi_1(S^1,1)\to\pi_1(X,A,1)\to\pi_0(A)\) 的两张映射。解释为什么末端 \(\pi_0(A)\) 的两点都能出现。
\end{exercise}
\begin{solution}
	把相对道路 \(\alpha\) 唯一提升成 \(\theta:I\to\mathbb R\)，使 \(\theta(0)=0\)、\(\alpha(t)=e^{2\pi\mathrm{i}\theta(t)}\)。终点条件等价于 \(\theta(1)\in\frac12\mathbb Z\)。在相对同伦中，终点连续地留在离散集 \(A\) 内，因而恒定；提升后的终值也连续地留在离散集 \(\frac12\mathbb Z\) 内，故不变。反过来，两条提升若终值相同，其线性插值固定两端，投影后给出相对同伦。因此 \([\alpha]\mapsto2\theta(1)\) 是到 \(\mathbb Z\) 的尖点集双射，指定元素对应 \(0\)。

	绕数为 \(k\) 的闭环路具有提升终值 \(k\)，所以第一张映射为 \(k\mapsto2k\)；第二张映射把偶数送到 \(\{1\}\)、奇数送到 \(\{-1\}\)。两点在 \(S^1\) 中属于同一道路分支，分别由常值道路与半圆道路到达，所以都出现。这里 \(\mathbb Z\) 用于标记尖点集；原来的相对道路没有自动给定以这个标号相加为定义的拼接运算。
\end{solution}

\begin{exercise}\label{g2:ex:topological-group-simple-action}
	设 \(G\) 为拓扑群，单位元为 \(e\)。证明 \(\pi_1(G,e)\) 在 \(\pi_n(G,e)\) 上的基点运输作用对每个 \(n\ge1\) 都是平凡的，并由 \(n=1\) 的情形推出 \(\pi_1(G,e)\) 交换。
\end{exercise}
\begin{solution}
	设 \(f:(D^n,S^{n-1})\to(G,e)\)，\(\gamma:I\to G\) 是基点环路。点乘公式
	\[
		H(z,s)=\gamma(s)f(z)
	\]
	定义连续同伦，边界值为 \(\gamma(s)\)，初映射和终映射都为 \(f\)。\cref{g2:prop:basepoint-action}证明中的移动边界刻画给出 \(c_\gamma[f]=[H(-,1)]=[f]\)，故作用平凡。对 \(n=1\)，这个作用就是 \(a\cdot g=g^{-1}ag\)，于是所有 \(a,g\in\pi_1(G,e)\) 都满足 \(g^{-1}ag=a\)，即 \(ag=ga\)。证明只使用群乘法连续，无需假设 \(G\) 本身交换。
\end{solution}

\begin{exercise}\label{g2:ex:fibration-composition}
	设 \(p:E\rightarrow B\)、\(q:B\rightarrow C\) 都是 Serre 纤维化。证明 \(qp\) 也是 Serre 纤维化，并证明对 Hurewicz 纤维化也有相同结论。若 \(c\in C\)，说明 \(p:(qp)^{-1}(c)\rightarrow q^{-1}(c)\) 为什么仍为相应类型的纤维化。
\end{exercise}
\begin{solution}
	设允许的参数空间为 \(Z\)，底同伦为 \(H:Z\times I\rightarrow C\)，初始提升为 \(a:Z\rightarrow E\)，满足 \(qpa=H|_{Z\times\{0\}}\)。先由 \(q\) 的提升性质，取以 \(pa\) 为初值的 \(H_B:Z\times I\rightarrow B\)，使 \(qH_B=H\)。再由 \(p\) 的提升性质，取以 \(a\) 为初值的 \(H_E\)，使 \(pH_E=H_B\)。于是 \(qpH_E=H\)，且初值正确。取 \(Z=D^m\) 就得到 Serre 情形，允许任意 \(Z\) 就得到 Hurewicz 情形。

	最后的映射是 \(p\) 沿包含 \(q^{-1}(c)\hookrightarrow B\) 的拉回：拉回空间中的配对 \((b,e)\) 满足 \(b=p(e)\) 和 \(q(b)=c\)，所以与 \((qp)^{-1}(c)\) 同胚。由\cref{g2:prop:fibration-pullback}即得结论。
\end{solution}

\begin{exercise}\label{g2:ex:homotopic-maps-homotopy-fibers}
	设 \(f_0,f_1:X\rightarrow B\) 由连续同伦 \(H:X\times I\rightarrow B\) 连接，固定 \(b_0\in B\)。构造 \(\operatorname{hofib}_{b_0}(f_0)\) 与 \(\operatorname{hofib}_{b_0}(f_1)\) 之间的同伦等价。说明所构造映射怎样作用于路径坐标，并证明两个复合确实同伦于恒等映射。
\end{exercise}
\begin{solution}
	令 \(h_x(s)=H(x,s)\)，它从 \(f_0(x)\) 走到 \(f_1(x)\)。道路拼接采用先左后右的约定，故定义
	\[
		\Phi(x,\gamma)=(x,\bar h_x*\gamma),\qquad
		\Psi(x,\eta)=(x,h_x*\eta).
	\]
	例如 \(\bar h_x*\gamma\) 从 \(f_1(x)\) 经 \(f_0(x)\) 到 \(b_0\)，所以 \(\Phi\) 的靶空间正确。紧开拓扑的指数律与道路的分段拼接公式保证二者连续。

	复合 \(\Psi\Phi\) 的路径坐标是 \(h_x*(\bar h_x*\gamma)\)。先连续移动两个分点，将三段道路的参数分配从 \((1/2,1/4,1/4)\) 改为 \((1/4,1/4,1/2)\)，便得到 \((h_x*\bar h_x)*\gamma\)。这一步只改变走路的速度，端点保持不变。对 \(r\in I\)，在前半段使用往返道路
	\[
		\ell_{x,r}(t)=
		\begin{cases}
			h_x(2(1-r)t),&0\le t\le1/2,\\
			h_x(2(1-r)(1-t)),&1/2\le t\le1.
		\end{cases}
	\]
	当 \(r=0\) 时它是 \(h_x*\bar h_x\)，当 \(r=1\) 时它是 \(f_0(x)\) 处的常值道路。于是 \(\ell_{x,r}*\gamma\) 给出定端点的连续变形；最后将常值段所占时间从 \(1/2\) 缩到零，就得到 \(\gamma\)。每一步都保持 \(x\)，并可由伴随公式验证联合连续。因此 \(\Psi\Phi\simeq\operatorname{id}\)。对 \(\Phi\Psi\)，用同样公式缩去 \(\bar h_x*h_x\) 即得 \(\Phi\Psi\simeq\operatorname{id}\)。
\end{solution}

\begin{exercise}\label{g2:ex:square-map-homotopy-fiber}
	令 \(f:\mathbb R\rightarrow\mathbb R\)、\(f(x)=x^2\)。分别求 \(f^{-1}(-1)\) 与 \(f^{-1}(1)\)，并证明对任意 \(b\in\mathbb R\)，\(\operatorname{hofib}_b(f)\) 都可缩。用一个道路提升问题证明 \(f\) 不是 Serre 纤维化。
\end{exercise}
\begin{solution}
	通常纤维分别为空集和两点集 \(\{-1,1\}\)。同伦纤维中的点则为 \((x,\gamma)\)，其中 \(\gamma(0)=x^2\)、\(\gamma(1)=b\)。先把路径线性变到连接两端的线段：
	\[
		\gamma_s(t)=(1-s)\gamma(t)+s\bigl((1-t)x^2+tb\bigr).
	\]
	这保持端点，且联合连续，因而将同伦纤维强形变收缩到由线段路径组成的子空间。该子空间由 \(x\) 唯一参数化，与 \(\mathbb R\) 同胚。再令 \(x_s=(1-s)x\)，并使用路径 \(t\mapsto(1-t)x_s^2+tb\)，就把它缩到 \(x=0\) 对应的一点。因此同伦纤维总是非空且可缩。

	取底道路 \(H(t)=-t\)，并指定在时刻零提升到 \(x=0\)。若存在提升 \(\widetilde H\)，则 \(\widetilde H(t)^2=-t\)，对任意 \(t>0\) 均不可能。\(f\) 连 \(D^0\) 参数的同伦提升性质都不满足，故不是 Serre 纤维化。
\end{solution}

\begin{exercise}\label{g2:ex:fibration-section-splitting}
	设 \(p:E\to B\) 是 Serre 纤维化，\(B\) 与基点纤维 \(F=p^{-1}(b_0)\) 都道路连通，且有截面 \(s:B\to E\)，取 \(e_0=s(b_0)\)。证明对 \(n\ge2\)，有 \(\pi_n(E,e_0)\cong\pi_n(F,e_0)\oplus\pi_n(B,b_0)\)；对 \(n=1\)，相应结论是分裂群扩张。用 Klein 瓶到圆周的投影说明基本群不必是直积。
\end{exercise}
\begin{solution}
	因为 \(p_*s_*=\operatorname{id}\)，每个 \(p_*\) 都满射。长正合列在底空间群处的正合性给出 \(\partial_n=0\)，于是对每个 \(n\ge1\)，都有分裂短正合列
	\[
		1\longrightarrow\pi_n(F)\xrightarrow{i_*}\pi_n(E)\xrightarrow{p_*}\pi_n(B)\longrightarrow1.
	\]
	在交换群次数也可把两端单位群写成 \(0\)。其中 \(i_*\) 的单射还使用了前一连接同态 \(\partial_{n+1}=0\)。对于 \(n\ge2\)，两个子群均交换，映射 \((a,b)\mapsto i_*a+s_*b\) 给出直和同构；任一 \(x\) 唯一写成 \((x-s_*p_*x)+s_*p_*x\)，前一项来自纤维，这证明满射及唯一性。对于 \(n=1\)，仍有唯一分解 \(x=i_*a\,s_*b\)，但 \(s_*b\) 对 \(i_*\pi_1(F)\) 的共轭作用可以非平凡，故得到半直积。

	具体取
	\[
		K=(S^1\times I)/((z,1)\sim(\bar z,0)),\qquad p([z,t])=e^{2\pi\mathrm i t}.
	\]
	在圆周的开弧上用时间坐标作平凡化，跨越接缝时纤维坐标变为 \(\bar z\)，故 \(p\) 是圆周丛；\(s(e^{2\pi\mathrm i t})=[1,t]\) 是截面。记纤维正向环路为 \(a\)，截面环路为 \(b\)。沿时间方向把纤维环路推到接缝，\(z\mapsto\bar z\) 将正向绕行变为负向绕行，所以
	\[
		b^{-1}ab=a^{-1}.
	\]
	上面的分裂短正合列表明 \(a\) 生成一个嵌入的 \(\mathbb Z\)，每个元素唯一写成 \(a^m b^n\)，故
	\[
		\pi_1(K)\cong\langle a,b\mid b^{-1}ab=a^{-1}\rangle\cong\mathbb Z\rtimes\mathbb Z.
	\]
	这里 \(a\) 有无穷阶，\(a\ne a^{-1}\)，所以基本群非交换，与 \(\mathbb Z^2\) 不同。
\end{solution}

\begin{exercise}\label{g2:ex:weak-equivalence-homotopy-fiber}
	设 \(f:(X,x_0)\to(Y,y_0)\) 为道路连通空间之间的带基点映射。证明 \(f\) 是弱同伦等价，当且仅当 \(F_f=\operatorname{hofib}_{y_0}(f)\) 道路连通且所有正次同伦群均为零。解释证明中为什么必须单独处理 \(\pi_0(F_f)\)。
\end{exercise}
\begin{solution}
	以\cref{g2:prop:fibration-replacement}将 \(f\) 写成 \(X\xrightarrow{j}E_f\xrightarrow{p_f}Y\)。\(j\) 是保持基点的同伦等价，所以 \(E_f\) 道路连通，且 \((p_f)_*\) 与 \(f_*\) 经 \(j_*\) 对应。若 \(f_*\) 在所有正次同伦群上均为同构，对 \(n\ge1\) 考察片段
	\[
		\pi_{n+1}(E_f)\longrightarrow\pi_{n+1}(Y)\longrightarrow\pi_n(F_f)
		\longrightarrow\pi_n(E_f)\longrightarrow\pi_n(Y).
	\]
	第一个映射满射，使进入 \(\pi_n(F_f)\) 的映射为零；最后一个映射单射，使 \(\pi_n(F_f)\to\pi_n(E_f)\) 的像为零。正合性于是迫使 \(\pi_n(F_f)=0\)。

	此外，\(E_f\) 道路连通，故\cref{g2:prop:fibration-component-action}说明 \(\pi_1(Y)\) 在 \(\pi_0(F_f)\) 上传递作用。指定分支的稳定子是 \(\operatorname{im}(p_f)_*\)，而 \(\pi_1(X)\to\pi_1(Y)\) 满射使这个稳定子等于整个群。因此指定分支的轨道只有一点，传递性又说明全部 \(\pi_0(F_f)\) 只有这一点。基点变化随即使纤维内任何基点的正次同伦群都为零。

	反过来，若 \(F_f\) 道路连通且正次同伦群全为零，则对 \(n\ge2\)，长正合列直接给出 \(\pi_n(E_f)\cong\pi_n(Y)\)。在一次处，\(\pi_1(F_f)=0\) 保证单射，\(\pi_0(F_f)\) 只有一点保证满射，故也为同构。源与靶都道路连通，零次条件自动满足，遂得 \(f\) 为弱同伦等价。单独检查零次项是必要的：所有正次群均为零的空间仍可能有多个道路分支。
\end{solution}

\begin{exercise}\label{g2:ex:hopf-pullback-degree}
	设 \(g:(S^2,b)\to(S^2,b)\) 是度为 \(d\in\mathbb Z\) 的带基点连续映射。沿 \(g\) 拉回复 Hopf 丛 \(h:S^3\to S^2\)，得到
	\[
		E_d=\{(x,z)\in S^2\times S^3:g(x)=h(z)\},\qquad p_d(x,z)=x.
	\]
	证明 \(E_d\) 道路连通，计算 \(\pi_1(E_d)\)、\(\pi_2(E_d)\)，并比较 \(\pi_k(E_d)\) 与 \(\pi_k(S^2)\)（\(k\ge3\)）。
\end{exercise}
\begin{solution}
	\cref{g2:prop:fibration-pullback}给出纤维为 \(S^1\) 的局部平凡丛。底空间与纤维均道路连通，先提升底空间道路，再在终点纤维内连接，便说明总空间道路连通。

	先确认度与二次同伦群的关系。由\cref{g2:prop:complex-hopf-homotopy}，\(\pi_2(S^2)\cong\mathbb Z\)；由\cref{g2:lem:degree-homotopy-additivity}，度是到 \(\mathbb Z\) 的群同态，而且恒等映射的度为 \(1\)，所以它满射。无限循环群到自身的满同态是同构，故 \(\deg:\pi_2(S^2)\to\mathbb Z\) 是同构。\cref{g2:prop:degree-basic}的复合公式因而说明 \(g_*\) 在此识别下是乘 \(d\)。

	拉回方块在纤维上为恒等，自然性给出 \(\partial^{(d)}_2=\partial^{\mathrm H}_2g_*\)。Hopf 丛的 \(\partial^{\mathrm H}_2\) 是无限循环群之间的同构，故 \(\partial^{(d)}_2\) 的核、像与乘 \(d\) 的核、像相同。长正合列成为
	\[
		0\longrightarrow\pi_2(E_d)\longrightarrow\mathbb Z\xrightarrow{\ \pm d\ }\mathbb Z
		\longrightarrow\pi_1(E_d)\longrightarrow0.
	\]
	正负号取决于顶面约定与球面、圆周生成元的对应，不影响核与像。因此
	\[
		\pi_1(E_d)\cong\mathbb Z/d\mathbb Z,\qquad
		\pi_2(E_d)\cong
		\begin{cases}\mathbb Z,&d=0,\\0,&d\ne0.\end{cases}
	\]
	这里约定 \(\mathbb Z/0\mathbb Z=\mathbb Z\)。当 \(k\ge3\) 时，\(\pi_k(S^1)=\pi_{k-1}(S^1)=0\)，所以 \((p_d)_*:\pi_k(E_d)\to\pi_k(S^2)\) 为同构。底映射的度改变低阶基本群，却不改变这个高阶比较。
\end{solution}

\begin{exercise}\label{g2:ex:hopf-transitions-sections}
	在\cref{g2:prop:hopf-bundles}的两张坐标图上，令 \(w=q_1q_0^{-1}\)、\(v=q_0q_1^{-1}\)，并沿用局部截面 \(s_0,s_1\)。求交集上的右过渡函数 \(g(w)\)，使 \(s_1(w^{-1})=s_0(w)g(w)\)。再证明复、四元数两种 Hopf 丛均无连续整体截面。
\end{exercise}
\begin{solution}
	交集上 \(w\ne0\)。把分母同乘 \(|w|\)，得到
	\[
		s_1(w^{-1})=\frac{(|w|w^{-1},|w|)}{\sqrt{1+|w|^2}}
		=s_0(w)\frac{\bar w}{|w|}.
	\]
	故 \(g(w)=\bar w/|w|\)。四元数情形中右侧第二坐标的核验是 \(w(\bar w/|w|)=|w|\)，不能任意颠倒系数。复情形在 \(|w|=1\) 上得到 \(g(e^{\mathrm{i}\theta})=e^{-\mathrm{i}\theta}\)。

	若 \(s:S^d\rightarrow S^{2d-1}\) 是 \(h_{\mathbb K}\) 的截面，则 \(h_{\mathbb K}s=\operatorname{id}_{S^d}\)。给总空间选基点 \(s(b)\)，这给出 \(h_{\mathbb K*}s_*\) 为 \(\pi_d(S^d,b)\) 上的恒等同态。但是 \(d<2d-1\)，故 \(\pi_d(S^{2d-1})=0\)。于是 \(\pi_d(S^d)\) 必为零，与恒等映射的度为 \(1\) 矛盾。因此两种丛都无连续整体截面，特别地都不是乘积丛。
\end{solution}

\begin{exercise}\label{g2:ex:projective-higher-homotopy}
	设整数 \(m\ge1\)。从单位向量到其复直线的投影 \(p:S^{2m+1}\rightarrow\mathbb{CP}^m\) 出发，证明
	\[
		\pi_1(\mathbb{CP}^m)=0,\qquad\pi_2(\mathbb{CP}^m)\cong\mathbb Z,
	\]
	以及 \(p_*:\pi_k(S^{2m+1})\rightarrow\pi_k(\mathbb{CP}^m)\) 在 \(k\ge3\) 时为同构。由此确定 \(3\le k\le2m\) 时的群，并说明标准包含 \(\mathbb{CP}^1\hookrightarrow\mathbb{CP}^m\) 在 \(\pi_2\) 上的作用。
\end{exercise}
\begin{solution}
	在第 \(j\) 个齐次坐标非零的图上，把该坐标归一为正实数，得到连续单位代表。于是
	\[
		p^{-1}(U_j)\longrightarrow U_j\times S^1,\qquad
		z\longmapsto([z],z_j/|z_j|)
	\]
	为平凡化，逆映射将所选单位代表乘以第二个坐标。因此 \(p\) 是纤维为 \(S^1\) 的局部平凡丛，且底空间是连通球面的连续像，因而道路连通。

	\(S^{2m+1}\) 的 \(\pi_1,\pi_2\) 都为零，\(S^1\) 道路连通；长正合列给出 \(\pi_1(\mathbb{CP}^m)=0\) 以及连接同构 \(\pi_2(\mathbb{CP}^m)\cong\pi_1(S^1)=\mathbb Z\)。在 \(k\ge3\) 的片段中，\(\pi_k(S^1)=\pi_{k-1}(S^1)=0\)，所以 \(p_*\) 为同构；当 \(3\le k\le2m\) 时，总空间的相应群又由胞腔逼近为零。

	把 \(\mathbb C^2\) 作为前两个坐标放入 \(\mathbb C^{m+1}\)，得到复 Hopf 丛到此圆周丛的映射，在基点纤维上是恒等映射。连接同态的自然性给出交换方块
	\[
		\begin{tikzcd}
			\pi_2(\mathbb{CP}^1) \arrow[r] \arrow[d,"\partial"',"\cong"] & \pi_2(\mathbb{CP}^m) \arrow[d,"\partial","\cong"']\\
			\pi_1(S^1) \arrow[r,"\operatorname{id}"'] & \pi_1(S^1).
		\end{tikzcd}
	\]
	因此标准包含在 \(\pi_2\) 上为同构，并保持由同一正向纤维环路指定的生成元。
\end{solution}

\begin{exercise}\label{g2:ex:two-frames-homotopy}
	设 \(n\ge3\)，令
	\[
		V_2(\mathbb R^n)=\{(u,v)\in S^{n-1}\times S^{n-1}:\langle u,v\rangle=0\}.
	\]
	证明取第一向量给出纤维为 \(S^{n-2}\) 的局部平凡丛，并计算 \(V_2(\mathbb R^n)\) 的 \(\pi_1\) 与 \(\pi_2\)。对 \(n=4\)，可把 \(\mathbb R^4\) 识别为 \(\mathbb H\)，考察 \((u,v)\mapsto(u,\bar u v)\)。
\end{exercise}
\begin{solution}
	在 \(u_0\) 附近连续选取正交矩阵 \(Q(u)\)，使第一列为 \(u\)，则 \(v=Q(u)(0,w)\) 唯一对应 \(w\in S^{n-2}\)。这给出所需平凡化。底与纤维均道路连通，利用道路提升再在终点纤维中连接，得到总空间道路连通。

	\(n=3\) 时，映射 \((u,v)\mapsto(u,v,u\times v)\) 是到 \(\operatorname{SO}(3)\) 的同胚，逆映射取前两列。因此 \(\pi_1(V_2(\mathbb R^3))=\mathbb Z/2\mathbb Z\)，\(\pi_2(V_2(\mathbb R^3))=0\)。

	\(n=4\) 时，\(\operatorname{Re}(\bar u v)=\langle u,v\rangle=0\)，且 \(|\bar u v|=1\)。故
	\[
		V_2(\mathbb R^4)\longrightarrow S^3\times S^2,\qquad(u,v)\longmapsto(u,\bar u v)
	\]
	为同胚，逆映射为 \((u,w)\mapsto(u,uw)\)，其中 \(S^2\) 是单位纯虚四元数球面。映射到积空间及其同伦逐坐标确定，因而 \(\pi_j(X\times Y)\cong\pi_j(X)\times\pi_j(Y)\)。现在 \(\pi_1(S^3)=\pi_1(S^2)=\pi_2(S^3)=0\)，而 \(\pi_2(S^2)=\mathbb Z\)，得到 \(\pi_1(V_2(\mathbb R^4))=0\)、\(\pi_2(V_2(\mathbb R^4))=\mathbb Z\)。

	当 \(n\ge5\) 时，纤维 \(S^{n-2}\) 与底 \(S^{n-1}\) 的一、二次同伦群均为零。长正合列立即给出 \(\pi_1(V_2(\mathbb R^n))=\pi_2(V_2(\mathbb R^n))=0\)。
\end{solution}

\begin{exercise}\label{g2:ex:unitary-root-cover}
	设 \(n\ge1\)、\(m\ge2\) 为整数，定义
	\[
		P_m=\{(A,z)\in\operatorname{U}(n)\times S^1:\det A=z^m\},\qquad p(A,z)=A.
	\]
	证明 \(p\) 是连通的 \(m\) 重覆盖；计算以 \((I,1)\) 为基点的 \(\pi_1(P_m)\)、\(\pi_2(P_m)\)，并说明 \(p_*\pi_1(P_m)\) 在 \(\pi_1(\operatorname{U}(n))\cong\mathbb Z\) 中是哪一个子群。
\end{exercise}
\begin{solution}
	\(p\) 是覆盖 \(S^1\rightarrow S^1,z\mapsto z^m\) 沿 \(\det\) 的拉回。在行列式落入一条真开圆弧的邻域上，选定连续实角 \(\theta(A)\)，其 \(m\) 个互不相同的根为
	\[
		\exp\bigl(2\pi\mathrm{i}(\theta(A)+r)/m\bigr),\qquad 0\le r<m.
	\]
	这直接给出 \(m\) 个覆盖叶。

	\(\operatorname{U}(n)\) 道路连通，故任何 \((A,z)\) 都可沿一条底道路的提升连到 \(I\) 上某点。正文给出的基本群生成环路 \(\gamma(t)=\operatorname{diag}(1,\ldots,1,e^{2\pi\mathrm{i}t})\) 从 \((I,1)\) 提升为 \((\gamma(t),e^{2\pi\mathrm{i}t/m})\)，终点是 \((I,e^{2\pi\mathrm{i}/m})\)。重复提升可连到 \(I\) 上所有 \(m\) 个点，所以 \(P_m\) 道路连通。

	行列式绕数为 \(k\) 的环路，其指定提升的根坐标终值为 \(e^{2\pi\mathrm{i}k/m}\)。因此闭合提升恰对应 \(k\in m\mathbb Z\)。由覆盖基本群的提升判据，\(p_*\) 单射且像为 \(m\mathbb Z\)，所以 \(\pi_1(P_m)\cong m\mathbb Z\cong\mathbb Z\)。离散纤维使覆盖在二次同伦群上诱导同构，故 \(\pi_2(P_m)\cong\pi_2(\operatorname{U}(n))=0\)。
\end{solution}

\begin{exercise}\label{g2:ex:rp2-higher-action}
	利用对径覆盖 \(q:S^2\rightarrow\mathbb{RP}^2\)，证明 \(\pi_2(\mathbb{RP}^2)\cong\mathbb Z\)，并计算 \(\pi_1(\mathbb{RP}^2)\cong\mathbb Z/2\mathbb Z\) 对它的作用。
\end{exercise}
\begin{solution}
	小球面邻域可以选得与其对径邻域不交，故 \(q\) 为二重覆盖；\(S^2\) 单连通。因此底基本群由环路提升的两个可能终点决定，恰为 \(\mathbb Z/2\mathbb Z\)。长正合列及离散纤维给出 \(q_*:\pi_2(S^2)\rightarrow\pi_2(\mathbb{RP}^2)\) 的同构，故目标为 \(\mathbb Z\)。

	先把 \(\pi_2(S^2)\) 的整数具体识别为度。\cref{g2:lem:degree-homotopy-additivity}给出群同态 \(\deg:\pi_2(S^2)\rightarrow\mathbb Z\)，其在恒等映射上取 \(1\)，所以满射；源又已由 Hopf 丛证明为无限循环群，故此同态为同构。

	取球面基点 \(x_0\)。底空间非平凡环路从 \(x_0\) 出发提升后终止于 \(-x_0\)。把定义基本群作用的移动基点同伦提升到球面，就将球面代表 \(f\) 换为对径映射 \(a(x)=-x\) 与 \(f\) 的复合，并沿提升道路把基点识别回来。球面单连通保证这次基点识别只依赖端点；在普通同伦意义下移动基点不改变度。因此该作用在度上是乘以 \(\deg a\)。由\cref{g2:prop:reflection-antipodal-degree}，\(\deg a=(-1)^3=-1\)，所以非单位基本群元素把整数 \(k\) 送到 \(-k\)。
\end{solution}

